\chapter{絶対値と三角不等式}\label{ch:text1}

\begin{textref}{第1章 1.1節 定義1.1, 命題1.2, 練習問題1.2}
  定義1.1(絶対値), その直後に列挙された性質, 命題1.2(三角不等式)を扱う. テキストの命題1.2の証明は両辺を2乗する方法であるが, 本章では「$-c \leq x \leq c \to \abs{x} \leq c$」を経由する証明を形式化する.
\end{textref}

\section{絶対値の定義}

\begin{dfn}[テキスト定義1.1]\label{dfn:abs}
  $\varphi_{\abs\cdot}(x, y) :\equiv (0 \leq x \land y = x) \lor (\neg\, 0 \leq x \land y = -x)$ とおく. 任意の $x$ について $\varphi_{\abs\cdot}(x, x)$ または $\varphi_{\abs\cdot}(x, -x)$ が成り立ち(排中律), $\varphi_{\abs\cdot}(x, y) \land \varphi_{\abs\cdot}(x, y') \to y = y'$ も明らかなので, $\ZFC \vdash \forall x\,\exists ! y\,\varphi_{\abs\cdot}(x, y)$ である. よって関数記号 $\abs{x}$ を定義公理
  \[ \forall x\,\bigl((0 \leq x \land \abs{x} = x) \lor (\neg\, 0 \leq x \land \abs{x} = -x)\bigr) \]
  で導入する(第\ref{ch:definitions}章). $x \in \R$ のとき $\neg\, 0 \leq x$ は $x < 0$ と同値である([A2]). テキストの定義はこれと一致する.
\end{dfn}

\section{絶対値の基本性質}

\begin{lem}[テキスト定義1.1の直後の性質]\label{lem:abs-basic}
  以下で現れる変数はすべて $\R$ の元とする.
  \begin{enumerate}[label={[A\arabic*]}]
    \item $0 \leq x \to \abs{x} = x$ (集合一般).
    \item $\neg\, y \leq x \to x < y$.
    \item $x < 0 \to \abs{x} = -x$ (集合一般).
    \item $x \leq 0 \to \abs{x} = -x$.
    \item $\abs{x} \in \R$.
    \item $0 \leq \abs{x}$.
    \item $\abs{-x} = \abs{x}$.
    \item $-\abs{x} \leq x \land x \leq \abs{x}$.
    \item $\abs{x} \leq y \to -y \leq x \land x \leq y$.
    \item $-y \leq x \to x \leq y \to \abs{x} \leq y$.
    \item $\abs{x} < y \to -y < x \land x < y$.
    \item $-y < x \to x < y \to \abs{x} < y$.
    \item $\abs{x - y} = \abs{y - x}$.
    \item[{[A14]}] $x \neq 0 \to 0 < \abs{x}$.
    \item[{[A15]}] $x \neq y \to x - y \neq 0$.
    \item[{[A16]}] $0 \leq x \to 0 \leq y \to 0 \leq x y$.
    \item[{[A17]}] $\abs{x y} = \abs{x}\,\abs{y}$.
  \end{enumerate}
\end{lem}
\begin{proof}[通常の証明]
  いずれも $0 \leq x$ と $\neg\, 0 \leq x$ (すなわち $x < 0$)の場合分けである. 例えば [A8]: $0 \leq x$ なら $\abs x = x$ であり, $-x \leq 0 \leq x$. $x < 0$ なら $\abs x = -x$ で $-\abs x = x$, また $x < 0 < -x$. [A10]: $0 \leq x$ なら $\abs x = x \leq y$. $x < 0$ なら $-y \leq x$ から $-x \leq y$. [A17] は $x, y$ の符号による4通りの場合分けである.
\end{proof}

\trees{AbsOfNonneg,NotLeLt,AbsOfNeg,AbsNonpos,AbsR,AbsNonneg,AbsNeg,AbsBound,AbsLe,LeAbs,AbsLt,LtAbs,AbsSubComm,AbsPos,SubNeZero,MulNonneg,AbsMul}
\section{三角不等式}

\begin{thm}[三角不等式. テキスト命題1.2]\label{thm:triangle}
  $\ZFC \vdash \forall a\forall b\,(a \in \R \to b \in \R \to \abs{a + b} \leq \abs{a} + \abs{b})$.
\end{thm}
\begin{proof}[通常の証明]
  [A8] より $-\abs{a} \leq a \leq \abs{a}$, $-\abs{b} \leq b \leq \abs{b}$. 辺々加えて([F21]) $-\abs{a} + (-\abs{b}) \leq a + b \leq \abs{a} + \abs{b}$. 左辺は $-(\abs{a} + \abs{b})$ に等しい([F7]). よって [A10] で $y := \abs{a} + \abs{b}$ として $\abs{a + b} \leq \abs{a} + \abs{b}$.
\end{proof}

\trees{Triangle}
\begin{cor}[テキスト練習問題1.2]\label{cor:reverse-triangle}
  $a, b \in \R$ について $\abs{a} - \abs{b} \leq \abs{a + b}$.
\end{cor}
\begin{proof}
  $a = (a + b) + (-b)$ なので, 三角不等式と [A7] より $\abs{a} \leq \abs{a + b} + \abs{b}$. 両辺に $-\abs{b}$ を加える([F19]). (この系の証明木は今後追加する.)
\end{proof}
